% exploration/Double_Exp_Iterated_Ln.tex
% Session: Double-Exponential and Iterated-Logarithm Extensions
% Date: 2026-04-21

\documentclass[11pt]{article}
\usepackage{amsmath,amssymb,amsthm,booktabs,xcolor}
\usepackage[margin=1.2in]{geometry}

\newtheorem{theorem}{Theorem}
\newtheorem{lemma}[theorem]{Lemma}
\newtheorem{proposition}[theorem]{Proposition}
\newtheorem{conjecture}[theorem]{Conjecture}
\newtheorem{definition}[theorem]{Definition}
\newtheorem{remark}[theorem]{Remark}

\title{Double-Exponential and Iterated-Logarithm Extensions\\
{\large Towers, Primitives, and SuperBEST Impact}}
\author{Exploration Session --- EML Framework}
\date{2026-04-21}

\begin{document}
\maketitle

\begin{abstract}
Following the cheap-addition breakthrough ($\mathrm{add} = 2n$ for all reals),
we investigate whether double-exponential or iterated-logarithm structures
give rise to new \emph{primitives} that compress EML tree depth further.
The main results are negative in the ``new primitive'' direction (towers always
decompose into depth-$\geq 2$ F16 sub-trees), but positive in the
``new universal function'' direction: we identify the \emph{iterated EML tower}
$\mathrm{EML}^{(k)}$ as the canonical depth-$k$ family, and prove that the
\emph{Ackermann-like} growth of these towers creates an asymptotic separation
from polynomial-depth ELC trees.
We also give node-count tables for functions whose natural expression involves
towers, and identify two conjectures forced by the structure.
\end{abstract}

\tableofcontents
\bigskip

%% ─────────────────────────────────────────────────────────────────────────────
\section{Setup: What is a Tower?}
\label{sec:setup}
%% ─────────────────────────────────────────────────────────────────────────────

\begin{definition}[EML tower of order $k$]
\[
  \mathrm{EML}^{(1)}(x) \;=\; \exp(x), \quad
  \mathrm{EML}^{(k)}(x) \;=\; \exp\!\bigl(\mathrm{EML}^{(k-1)}(x)\bigr).
\]
Dually, $\mathrm{LML}^{(k)}(x) = \ln^{(k)}(x)$ (the $k$-fold iterated logarithm).
\end{definition}

\noindent Key questions:
\begin{enumerate}
  \item Can $\mathrm{EML}^{(k)}(x)$ or $\mathrm{LML}^{(k)}(x)$ be expressed
        as a \emph{depth-1} (single-node) F16 application? (Answer: No for $k \geq 2$.)
  \item Does any function in common use \emph{require} a tower structure,
        making tower-based operators cost-saving?
  \item What is the asymptotic depth needed to express
        $\exp_k(x)$ (the $k$-th iterated exponential) in an EML tree?
\end{enumerate}

%% ─────────────────────────────────────────────────────────────────────────────
\section{Double-Exponential: $\exp(\exp(x))$}
\label{sec:double-exp}
%% ─────────────────────────────────────────────────────────────────────────────

\subsection{Minimal EML Tree}

\begin{proposition}
$\exp(\exp(x))$ requires exactly 2 F16 nodes. It cannot be expressed as a
single F16 operator applied to raw input $x$.
\end{proposition}

\begin{proof}
\textit{Upper bound:} $\mathrm{EML}(\mathrm{EML}(x,1),1) = \exp(\exp(x) - 0) = \exp(\exp(x))$.
Two nodes.

\textit{Lower bound:} Every single F16 application $F(x,y)$ takes the form
$f(\exp(\varepsilon x), \ln(\delta y))$. Setting $y = 1$ gives
$f(\exp(\varepsilon x), 0)$ (since $\ln(1)=0$). This equals either
$\exp(\varepsilon x)$, $-\exp(\varepsilon x)$, $0$, or is undefined.
None equal $\exp(\exp(x))$ for all $x$ (as $\exp(\exp(x))$ grows faster
than $\exp(x)$ and $\exp(\exp(x)) > \exp(x)$ for $x > 0$).
For general $y$: $f(\exp(\varepsilon x), \ln(\delta y)) = \exp(\exp(x))$
would require $\ln(\delta y)$ to encode $\exp(x) - \exp(\varepsilon x)$ for
all $x$, which is impossible for any fixed $y$ or $y$ depending only on $x$
without an additional function application.
\end{proof}

\subsection{Domain Analysis}

$\exp(\exp(x))$ is defined for all $x \in \mathbb{R}$.
Growth: $\exp(\exp(x)) \sim \exp(e^x)$ --- doubly exponential.
This is \emph{not} in the ELC class at depth 1, but \emph{is} at depth 2.

\subsection{Functions Requiring $\exp(\exp(x))$}

Which standard functions need double-exponential depth?

\begin{center}
\begin{tabular}{lcc}
\toprule
Function & Depth-2 EML construction & Notes \\
\midrule
$\exp(\exp(x))$ & $\mathrm{EML}(\mathrm{EML}(x,1),1)$ & 2n \\
$\exp(e^x - 1)$ & $\mathrm{EML}(\mathrm{EML}(x,1)-1, 1)$... & 3n \\
Tetration ${}^2 e = e^e$ & constant; depth 2 tree & 2n (constant) \\
$\exp(x^2)$ & $= \exp(\mathrm{EPL}(2,x))$: 2n & 2n \\
$\exp(\sqrt{x})$ & $= \exp(\mathrm{EPL}(0.5,x))$: 2n & 2n \\
\bottomrule
\end{tabular}
\end{center}

\noindent \textbf{Key observation:} $\exp(x^n) = \mathrm{EML}(\mathrm{EPL}(n,x), 1)$ at 2n.
This means all \emph{sub-Gaussian} exponentials are depth 2, and only
\emph{iterated} exponentials (where the exponent itself grows exponentially) genuinely
push to depth 3+.

%% ─────────────────────────────────────────────────────────────────────────────
\section{Iterated Logarithm: $\ln(\ln(y))$}
\label{sec:iterated-log}
%% ─────────────────────────────────────────────────────────────────────────────

\begin{proposition}
$\ln(\ln(y))$ (defined for $y > 1$) requires exactly 2 F16 nodes.
\end{proposition}

\begin{proof}
Upper bound: $\mathrm{EXL}(0, \mathrm{EXL}(0, y)) = \ln(\ln(y))$, 2 nodes.
Lower bound: same argument as double-exp; no single F16 evaluation
$f(\exp(\varepsilon x), \ln(\delta y))$ with constant first argument equals $\ln(\ln(y))$
for all admissible $y$.
\end{proof}

\subsection{Domain: $y > 1$ (so $\ln y > 0$)}

$\ln(\ln(y))$ is defined and real only for $y > 1$.
The iterated-log tower $\ln^{(k)}(y)$ has progressively restricted domain:
$y > 1$, $y > e$, $y > e^e$, \ldots

\subsection{Functions Requiring $\ln(\ln(y))$}

\begin{center}
\begin{tabular}{lcc}
\toprule
Function & Construction & Cost \\
\midrule
$\ln(\ln(y))$ & $\mathrm{EXL}(0,\mathrm{EXL}(0,y))$ & 2n \\
$\ln(\ln(y)) - \ln(\ln(z))$ & $\mathrm{EXL}(0,\mathrm{LLS}(y,z))$... & complex \\
$\ln\ln(e^e) = 1$ & constant & 0n \\
Ramanujan--Hardy: $\ln\ln$ in prime gaps & requires depth $\geq 2$ & 2n per $\ln\ln$ \\
\bottomrule
\end{tabular}
\end{center}

%% ─────────────────────────────────────────────────────────────────────────────
\section{Mixed Tower: $\ln(\exp(x) - \ln(y))$}
\label{sec:mixed}
%% ─────────────────────────────────────────────────────────────────────────────

This is $\mathrm{EXL}(0, \mathrm{EML}(x,y)) = \ln(\exp(x) - \ln(y))$,
a 2-node EML tree. It represents the \emph{log of EML}.

\begin{proposition}
$\ln(\exp(x) - \ln(y))$ (domain: $\exp(x) > \ln(y)$) is expressible in exactly
2 F16 nodes and equals $\mathrm{EXL}(0, \mathrm{EML}(x,y))$.
\end{proposition}

\noindent \textbf{Why interesting:} Composing $\ln$ with EML produces a
function in the \emph{inverse} EML class. These arise in:
\begin{itemize}
  \item Solution of $\exp(x) - \ln(y) = e^z$: gives $z = \ln(\mathrm{EML}(x,y))$.
  \item Double-application of the EML gate in autofeedback loops.
  \item The ``phantom orbit'' analysis: iterating $\ln \circ \mathrm{EML}$ is
        structurally different from iterating EML itself.
\end{itemize}

%% ─────────────────────────────────────────────────────────────────────────────
\section{Asymptotic Depth: Tower Hierarchy Separation}
\label{sec:depth}
%% ─────────────────────────────────────────────────────────────────────────────

\begin{definition}[ELC depth-$k$ class]
$\mathrm{ELC}_k(\mathbb{R})$ is the set of functions expressible by EML trees
of depth $\leq k$.
\end{definition}

\begin{theorem}[Strict tower separation]
\label{thm:tower-sep}
For each $k \geq 0$, the function $\mathrm{EML}^{(k+1)}(x) = \exp_{k+1}(x)$
(the $(k+1)$-fold iterated exponential) lies in $\mathrm{ELC}_{k+1}$ but not
in $\mathrm{ELC}_k$.
\end{theorem}

\begin{proof}
\textit{Upper bound:} $\exp_{k+1}(x) = \mathrm{EML}(\exp_k(x), 1)$,
and $\exp_k(x) \in \mathrm{ELC}_k$ by induction. So $\exp_{k+1} \in \mathrm{ELC}_{k+1}$.

\textit{Lower bound:} Every function in $\mathrm{ELC}_k$ has growth bounded by
an $(k)$-fold iterated exponential (formal: the Liouville--Osgood tower bound
for iterated exp--log compositions). But $\exp_{k+1}(x)$ grows strictly faster
than any $\mathrm{ELC}_k$ function for large $x$ (by induction on $k$, using
the fact that $\mathrm{EML}$ at depth $k$ produces at most $k$-fold composed
exponentials).
\end{proof}

\begin{remark}
Theorem~\ref{thm:tower-sep} is the analogue of the \emph{Grzegorczyk hierarchy}
for elementary arithmetic functions, but in the EML/exp--log setting.
It confirms that the depth hierarchy $\mathrm{ELC}_0 \subsetneq \mathrm{ELC}_1
\subsetneq \cdots$ is strict at \emph{every} level, not just at the first few.
\end{remark}

%% ─────────────────────────────────────────────────────────────────────────────
\section{Are Tower-Based Operators Ever Primitives?}
\label{sec:primitives}
%% ─────────────────────────────────────────────────────────────────────────────

\begin{proposition}
No function of the form $f(\exp(\exp(x)), g(y))$ or $f(\exp(x), \ln(\ln(y)))$
can be a depth-1 ``primitive'' in any extension of F16 without inflating the
\emph{gate complexity} (i.e., encoding a depth-$\geq 2$ sub-tree in a single gate).
\end{proposition}

\begin{proof}
A depth-1 primitive $P(x,y)$ requires $P$ to be computable from $x$ and $y$
using exactly one application of a fixed function $f$. If $f = f(\exp(\exp(\cdot)), \cdot)$,
then evaluating $P$ at a machine node requires computing $\exp(\exp(x))$ as a
sub-computation---which requires 2 sequential exponential evaluations.
This violates the unit-cost assumption of the single-node model.
The gate complexity of $P$ is $\geq 2$.
\end{proof}

\noindent \textbf{Conclusion:} Double-exponential and iterated-logarithm forms are
\emph{not} primitives in any physically motivated single-node model.
They are depth-$k$ compositions of F16 nodes. The F16 basis is not extended by towers.

%% ─────────────────────────────────────────────────────────────────────────────
\section{The Iterated EML Map: Fixed-Point Tower}
\label{sec:fixedpoint}
%% ─────────────────────────────────────────────────────────────────────────────

Define the EML self-map $T : \mathbb{R} \to \mathbb{R}$ by
$T(x) = \mathrm{EML}(x,x) = \exp(x) - \ln(x)$ for $x > 0$.

\begin{proposition}
$T$ has a unique fixed point $x^* \approx 3.1696$ (the phantom attractor).
The iterates $T^{(k)}(x)$ converge to $x^*$ for all $x > 0$ in a neighbourhood
of $x^*$, and diverge to $+\infty$ for $x$ sufficiently small or large.
\end{proposition}

\noindent The tower $T^{(k)}(x)$ thus defines a natural depth-$k$ family,
and the phantom attractor is the ``tower fixed point'' of the EML gate.
This is structurally distinct from the double-exp tower $\exp_k(x)$:
the EML tower \emph{converges}, while the exp tower \emph{diverges}.

\begin{conjecture}[CONJ\_TOWER\_BASIN]
The basin of attraction of $x^*$ under $T(x) = \exp(x) - \ln(x)$ is
the interval $(0, \exp(x^*))$ (approximately $(0, 23.83)$), and for all
$x$ in this interval, $T^{(k)}(x) \to x^*$ monotonically for $k \geq k_0(x)$.
\end{conjecture}

%% ─────────────────────────────────────────────────────────────────────────────
\section{Node Count Table: Functions Involving Towers}
\label{sec:table}
%% ─────────────────────────────────────────────────────────────────────────────

\begin{center}
\begin{tabular}{lccp{5cm}}
\toprule
Function & EML tree depth & Node count & Construction \\
\midrule
$\exp(x)$        & 1 & 1n & $\mathrm{EML}(x,1)$ \\
$\exp(\exp(x))$  & 2 & 2n & $\mathrm{EML}(\mathrm{EML}(x,1),1)$ \\
$\exp_3(x)$      & 3 & 3n & iterating above \\
$\exp_k(x)$      & k & kn & by induction \\
$\ln(x)$         & 1 & 1n & $\mathrm{EXL}(0,x)$ \\
$\ln(\ln(x))$    & 2 & 2n & $\mathrm{EXL}(0,\mathrm{EXL}(0,x))$ \\
$\ln_k(x)$       & k & kn & by induction \\
$\exp(\ln(x))=x$ & 2 & 2n & but simplifies to $x$ (0n) \\
$\ln(\exp(x))=x$ & 2 & 2n & but simplifies to $x$ (0n) \\
$\exp(x)-\ln(x)$ (EML self) & 1 & 1n & $\mathrm{EML}(x,x)$ \\
$T^{(2)}(x)$     & 2 & 2n & $\mathrm{EML}(\mathrm{EML}(x,x),\mathrm{EML}(x,x))$... \\
\bottomrule
\end{tabular}
\end{center}

%% ─────────────────────────────────────────────────────────────────────────────
\section{Summary and Structural Comparison}
\label{sec:summary}
%% ─────────────────────────────────────────────────────────────────────────────

\begin{enumerate}
  \item \textbf{Double-exp and iterated-log are not new primitives.}
        They are depth-$k$ F16 trees. Adding them as gates inflates gate complexity.
  \item \textbf{The depth hierarchy is strict at every level} (Theorem~\ref{thm:tower-sep}).
        This is the EML analogue of the Grzegorczyk arithmetic hierarchy.
  \item \textbf{The EML self-map tower $T^{(k)}$ converges} to the phantom attractor
        $x^* \approx 3.1696$, distinguishing it from the divergent exp tower.
  \item \textbf{For SuperBEST:} tower-involving functions cost exactly $k$ nodes
        per $k$-fold nesting. No compression is possible beyond the single-gate level.
  \item \textbf{The only ``new primitives'' that survive scrutiny} are those in the
        EE/LL families (two-exp or two-log arguments), not tower forms.
\end{enumerate}

\begin{conjecture}[CONJ\_GRZEGORCZYK\_EML]
The ELC depth hierarchy $\mathrm{ELC}_0 \subsetneq \mathrm{ELC}_1 \subsetneq \cdots$
is in precise correspondence with the Grzegorczyk hierarchy
$\mathcal{E}^0 \subsetneq \mathcal{E}^1 \subsetneq \cdots$ restricted to
real-analytic functions, where the correspondence is given by:
$\mathrm{ELC}_k$ contains exactly those elementary functions whose
\emph{iterated exponential growth order} is $\leq k$.
\end{conjecture}

\end{document}
